
\documentclass{article}
\usepackage{colortbl}
\usepackage{makecell}
\usepackage{multirow}
\usepackage{supertabular}

\begin{document}

\newcounter{utterance}

\centering \large Interaction Transcript for game `codenames', experiment `unambiguous', episode 6 with reflexion{-}Qwen3.5{-}27B.
\vspace{24pt}

{ \footnotesize  \setcounter{utterance}{1}
\setlength{\tabcolsep}{0pt}
\begin{supertabular}{c@{$\;$}|p{.15\linewidth}@{}p{.15\linewidth}p{.15\linewidth}p{.15\linewidth}p{.15\linewidth}p{.15\linewidth}}
    \# & $\;$A & \multicolumn{4}{c}{Game Master} & $\;\:$B\\
    \hline

    \theutterance \stepcounter{utterance}  
    & & & \multicolumn{2}{p{0.3\linewidth}}{
        \cellcolor[rgb]{0.9,0.9,0.9}{
            \makecell[{{p{\linewidth}}}]{
                \texttt{\tiny{[GM$|$GM]}}
                {'hidden': {'team': ['toothpick'], 'innocent': ['nuisance', 'ozone', 'tequila', 'decorator', 'stewardess', 'motorcade', 'pond'], 'opponent': [], 'assassin': ['titanium']}, 'revealed': {'team': {'team': [], 'innocent': [], 'opponent': [], 'assassin': []}, 'opponent': {'team': [], 'innocent': [], 'opponent': [], 'assassin': []}}}
            }
        }
    }
    & & \\ \\

    \theutterance \stepcounter{utterance}  
    & & \multicolumn{4}{p{0.6\linewidth}}{
        \cellcolor[rgb]{0.9,0.9,0.9}{
            \makecell[{{p{\linewidth}}}]{
                \texttt{\tiny{[P1$\langle$GM]}}
                \texttt{You have attempted this same task before and did not succeed. These are the plans you wrote after those attempts.} \\
\texttt{Trial \#1: New plan: My previous strategy of giving clues for only one word at a time was inefficient, causing me to run out of turns before clearing my list. Specifically, I failed to group "toothpick" with other remaining words earlier in the game. For instance, "toothpick" and "ketchup" could have been grouped under a clue like "picnic" or "food" (if applicable to others), or "toothpick" and "jukebox" under "stick" (as in a drumstick or music stick, though weak). A better approach is to always scan the remaining list for a common semantic link that covers at least two words. In the final turn, instead of "wood" for just "toothpick," I should have looked for a broader category earlier. If I had grouped "toothpick" with "ketchup" (e.g., CLUE: "dip" or "snack" if "ketchup" was still there) or "toothpick" with "jukebox" (e.g., CLUE: "stick" if "jukebox" was still there), I would have cleared multiple words in one turn. The new strategy is to prioritize multi{-}word clues over single{-}word clues to maximize the number of words guessed per turn, ensuring I clear the list before the opponent does.} \\
\\ 
\texttt{Trial \#2: New plan: My previous attempt failed because I provided a clue ("wood") that was too specific to a single word ("toothpick"), resulting in an inefficient turn that likely allowed the opponent to win first. Although the current list only contains "toothpick," the strategy must be to always maximize the number of targets per clue. In future turns with multiple words, I will scan the list for a common semantic link (e.g., "dining" for "toothpick" and "ketchup") to group at least two words. If only one word remains, I will ensure the clue is the most direct and unambiguous association possible to guarantee the guess, rather than a generic material description that might be ambiguous or slow. The core rule is: never output a clue for a single word if a multi{-}word grouping is possible; if only one word remains, use the strongest direct association.} \\
\\ 
\texttt{Trial \#3: New plan: My previous attempt failed because I provided a clue ("stick") that was too generic and likely triggered a "distractor" word (e.g., "motorcade" or "decorator" might be associated with sticks in a stretch, or the partner simply guessed a distractor instead of the target). Although "toothpick" was the only target, the clue lacked specificity to distinguish it from the distractors. In future turns, even with a single word, I must choose a clue that is uniquely and strongly associated with the target while having zero semantic overlap with the distractor list. For "toothpick," instead of "stick," I should have used a more specific functional clue like "floss" (associated with dental hygiene) or "dental" to avoid ambiguity. The new strategy is to rigorously cross{-}reference potential clues against the distractor list to ensure the clue points *only* to the target words, prioritizing unique functional associations over generic material or shape descriptions.} \\
\\ 
\\ 
\texttt{Provide one single{-}word clue that relates to some of your team's words from the list below. You can choose to how many and to which words your clue relates to.} \\
\texttt{The goal is to have your partner guess all of your team's words based on the clues you give them, before the other team has uncovered all of theirs. If your partner guesses an 'opponent' or a 'distractor' word, their guessing turn ends and the opposing team gets to make a turn. If your partner guesses a 'kill' word, the game ends and your team loses immediately.} \\
\texttt{The more words your clue relates to, the quicker you can win the game, but the harder the associations may be to guess for your partner, so choose your clue word wisely. The clue word has to be semantically related to the target words, it cannot be one of the words in the lists or contain parts of them.} \\
\\ 
\texttt{Always give your single{-}word clue and your comma{-}separated list of related target words in the following format and make your answers as short as possible, never include any other text than is required in this form:} \\
\\ 
\texttt{CLUE: <WORD>} \\
\texttt{TARGETS: <TARGETS>} \\
\\ 
\texttt{Your team words are: toothpick.} \\
\texttt{Your opponent's team words are: .} \\
\texttt{Distractor words are: nuisance, ozone, tequila, decorator, stewardess, motorcade, pond.} \\
\texttt{The kill words are: titanium.} \\
            }
        }
    }
    & \\ \\

    \theutterance \stepcounter{utterance}  
    & \multicolumn{4}{p{0.6\linewidth}}{
        \cellcolor[rgb]{0.9,0.9,0.9}{
            \makecell[{{p{\linewidth}}}]{
                \texttt{\tiny{[P1$\rangle$GM]}}
                \texttt{CLUE: dental} \\
\texttt{TARGETS: toothpick} \\
            }
        }
    }
    & & \\ \\

    \theutterance \stepcounter{utterance}  
    & & & \multicolumn{2}{p{0.3\linewidth}}{
        \cellcolor[rgb]{0.9,0.9,0.9}{
            \makecell[{{p{\linewidth}}}]{
                \texttt{\tiny{[GM$|$GM]}}
                \texttt{dental} \\
            }
        }
    }
    & & \\ \\

    \theutterance \stepcounter{utterance}  
    & & & \multicolumn{2}{p{0.3\linewidth}}{
        \cellcolor[rgb]{0.9,0.9,0.9}{
            \makecell[{{p{\linewidth}}}]{
                \texttt{\tiny{[GM$|$GM]}}
                ['toothpick']
            }
        }
    }
    & & \\ \\

    \theutterance \stepcounter{utterance}  
    & & & \multicolumn{2}{p{0.3\linewidth}}{
        \cellcolor[rgb]{0.9,0.9,0.9}{
            \makecell[{{p{\linewidth}}}]{
                \texttt{\tiny{[GM$|$GM]}}
                {'word': 'toothpick', 'assignment': 'team'}
            }
        }
    }
    & & \\ \\

    \theutterance \stepcounter{utterance}  
    & & & \multicolumn{2}{p{0.3\linewidth}}{
        \cellcolor[rgb]{0.9,0.9,0.9}{
            \makecell[{{p{\linewidth}}}]{
                \texttt{\tiny{[GM$|$GM]}}
                \texttt{opponent has won} \\
            }
        }
    }
    & & \\ \\

    \theutterance \stepcounter{utterance}  
    & \multicolumn{4}{p{0.6\linewidth}}{
        \cellcolor[rgb]{0.95,0.95,0.95}{
            \makecell[{{p{\linewidth}}}]{
                \texttt{\tiny{[P2$\rightarrow$P1]}}
                \texttt{Post{-}episode Reflexion (not sent to GM):} \\
\texttt{New plan: My previous attempt failed because I provided a clue ("dental") that, while semantically related to "toothpick," was too broad and likely led my partner to guess a distractor word (such as "nuisance" or "decorator" in a stretch, or simply failing to guess the specific target quickly enough) or missed the target entirely, resulting in a loss. The core mistake was not rigorously filtering the clue against the specific distractor list to ensure a unique, high{-}probability association. In future turns, I will first list all potential clues for my target words and then explicitly cross{-}reference them against the distractor and kill word lists to eliminate any clue that could plausibly trigger a wrong guess. For "toothpick," instead of the generic "dental," I should have chosen a more specific functional or contextual clue like "floss" (strongly linked to dental hygiene but less likely to trigger "nuisance" or "ozone") or "cocktail" (if the context allows), ensuring the clue points exclusively to the target. The new strategy is: 1) Identify the strongest semantic link for the target(s). 2) Verify this link does not strongly associate with any distractor or kill word. 3) If only one word remains, prioritize the most specific, unambiguous association over a general category to minimize the risk of a distractor guess.} \\
            }
        }
    }
    & & \\ \\

\end{supertabular}
}

\end{document}
